%-----------------------------------------------------------------------------------------------------------------------------------------------%
%    The MIT License (MIT)
%
%    Copyright (c) 2015 Jan Küster
%
%    Permission is hereby granted, free of charge, to any person obtaining a copy
%    of this software and associated documentation files (the "Software"), to deal
%    in the Software without restriction, including without limitation the rights
%    to use, copy, modify, merge, publish, distribute, sublicense, and/or sell
%    copies of the Software, and to permit persons to whom the Software is
%    furnished to do so, subject to the following conditions:
%    
%    THE SOFTWARE IS PROVIDED "AS IS", WITHOUT WARRANTY OF ANY KIND, EXPRESS OR
%    IMPLIED, INCLUDING BUT NOT LIMITED TO THE WARRANTIES OF MERCHANTABILITY,
%    FITNESS FOR A PARTICULAR PURPOSE AND NONINFRINGEMENT. IN NO EVENT SHALL THE
%    AUTHORS OR COPYRIGHT HOLDERS BE LIABLE FOR ANY CLAIM, DAMAGES OR OTHER
%    LIABILITY, WHETHER IN AN ACTION OF CONTRACT, TORT OR OTHERWISE, ARISING FROM,
%    OUT OF OR IN CONNECTION WITH THE SOFTWARE OR THE USE OR OTHER DEALINGS IN
%    THE SOFTWARE.
%    
%
%-----------------------------------------------------------------------------------------------------------------------------------------------%


%============================================================================%
%
%    DOCUMENT DEFINITION
%
%============================================================================%

%we use article class because we want to fully customize the page and dont use a cv template
\documentclass[10pt,A4]{article}    


%----------------------------------------------------------------------------------------
%    ENCODING
%----------------------------------------------------------------------------------------

%we use utf8 since we want to build from any machine
\usepackage[utf8]{inputenc}        

%----------------------------------------------------------------------------------------
%    LOGIC
%----------------------------------------------------------------------------------------

% provides \isempty test
\usepackage{xifthen}

%----------------------------------------------------------------------------------------
%    FONT
%----------------------------------------------------------------------------------------

% some tex-live fonts - choose your own

%\usepackage[defaultsans]{droidsans}
%\usepackage[default]{comfortaa}
%\usepackage{cmbright}
\usepackage[default]{raleway}
%\usepackage{fetamont}
%\usepackage[default]{gillius}
%\usepackage[light,math]{iwona}
%\usepackage[thin]{roboto} 

% set font default
\renewcommand*\familydefault{\sfdefault}     
\usepackage[T1]{fontenc}

% more font size definitions
\usepackage{moresize}        


%----------------------------------------------------------------------------------------
%    PAGE LAYOUT  DEFINITIONS
%----------------------------------------------------------------------------------------

%debug page outer frames
%\usepackage{showframe}            


%define page styles using geometry
\usepackage[a4paper]{geometry}        

% for example, change the margins to 2 inches all round
\geometry{top=1.25cm, bottom=0cm, left=1.5cm, right=1.5cm}     

%use customized header
\usepackage{fancyhdr}                
\pagestyle{fancy}

%less space between header and content
\setlength{\headheight}{-5pt}        


%customize entries left, center and right
%\lhead{}
%\chead{}
%\rhead{}


%indentation is zero
\setlength{\parindent}{0mm}

%----------------------------------------------------------------------------------------
%    TABLE /ARRAY DEFINITIONS
%---------------------------------------------------------------------------------------- 

%for layouting tables
\usepackage{multicol}            
\usepackage{multirow}

%extended aligning of tabular cells
\usepackage{array}

\newcolumntype{x}[1]{%
>{\raggedleft\hspace{0pt}}p{#1}}%


%----------------------------------------------------------------------------------------
%    GRAPHICS DEFINITIONS
%---------------------------------------------------------------------------------------- 

%for header image
\usepackage{graphicx}

%for floating figures
\usepackage{wrapfig}
\usepackage{float}
%\floatstyle{boxed} 
%\restylefloat{figure}

%for drawing graphics        
\usepackage{tikz}                
\usetikzlibrary{shapes, backgrounds,mindmap, trees}


%----------------------------------------------------------------------------------------
%    Color DEFINITIONS
%---------------------------------------------------------------------------------------- 

\usepackage{color}

%accent color
\definecolor{sectcol}{RGB}{0,150,255}

%dark background color / accent color
\definecolor{bgcol}{RGB}{110,110,110}

%light background
\definecolor{softcol}{RGB}{225,225,225}

%----------------------------------------------------------------------------------------
%    HYPERLINKS
%----------------------------------------------------------------------------------------

\usepackage{hyperref}        
\hypersetup{
    colorlinks=true,
    linkcolor=sectcol,
    filecolor=magenta,
    urlcolor=sectcol,
    pdftitle={Overleaf Example},
    pdfpagemode=FullScreen,
    }


%============================================================================%
%
%
%    DEFINITIONS
%
%
%============================================================================%

%----------------------------------------------------------------------------------------
%     HEADER
%----------------------------------------------------------------------------------------

% remove top header line
\renewcommand{\headrulewidth}{0pt} 

%remove botttom header line
\renewcommand{\footrulewidth}{0pt}          

%remove pagenum
\renewcommand{\thepage}{}    

%remove section num        
\renewcommand{\thesection}{}            

%----------------------------------------------------------------------------------------
%     ARROW GRAPHICS in Tikz
%----------------------------------------------------------------------------------------

% a six pointed arrow poiting to the left
\newcommand{\tzlarrow}{(0,0) -- (0.2,0) -- (0.3,0.2) -- (0.2,0.4) -- (0,0.4) -- (0.1,0.2) -- cycle;}    

% include the left arrow into a tikz picture
% param1: fill color
%
\newcommand{\larrow}[1]
{\begin{tikzpicture}[scale=0.58]
     \filldraw[fill=#1!100,draw=#1!100!black]  \tzlarrow
 \end{tikzpicture}
}

% a six pointed arrow poiting to the right
\newcommand{\tzrarrow}{ (0,0.2) -- (0.1,0) -- (0.3,0) -- (0.2,0.2) -- (0.3,0.4) -- (0.1,0.4) -- cycle;}

% include the right arrow into a tikz picture
% param1: fill color
%
\newcommand{\rarrow}
{\begin{tikzpicture}[scale=0.7]
    \filldraw[fill=sectcol!100,draw=sectcol!100!black] \tzrarrow
 \end{tikzpicture}
}



%----------------------------------------------------------------------------------------
%    custom sections
%----------------------------------------------------------------------------------------

% create a coloured box with arrow and title as cv section headline
% param 1: section title
%
\newcommand{\cvsection}[1]
{
    \begin{center}
        \large\textcolor{sectcol}{\textbf{#1}}
    \end{center}
}

%create a coloured arrow with title as cv meta section section
% param 1: meta section title
%
\newcommand{\metasection}[2]
{
%\begin{tabular*}{1\textwidth}{r r}
\footnotesize{#2} \hspace*{\fill} \footnotesize{#1}\\[1pt]
%\end{tabular*}
}

%----------------------------------------------------------------------------------------
%     CV EVENT
%----------------------------------------------------------------------------------------

% creates a stretched box as cv entry headline followed by some paragraphs about 
% the work you did
% param 1:    event time i.e. 2014 or 2011-2014 etc.
% param 2:    event name (what did you do?)
% param 3:    institution (where did you work / study)
% param 4:    list of paragraphs
%
\newcommand{\cvevent}[4]
{

\begin{tabular*}{1\textwidth}{p{13.0cm}  x{4.5cm}}
    \textbf{#2} - \textcolor{bgcol}{#3} &   \vspace{2.5pt}\textcolor{sectcol}{#1}
\end{tabular*}

\vspace{-8pt}
\textcolor{softcol}{\hrule}
\vspace{6pt}

    \foreach \desc in {#4}{
        $\cdot$ \desc\\[3pt]
    }

\vspace{3pt}
}

% creates a stretched box as 
\newcommand{\cveventmeta}[2]
{
    \mbox{\mystrut \hspace{87pt}\textit{#1}}\\
    #2
}

%----------------------------------------------------------------------------------------
% CUSTOM STRUT FOR EMPTY BOXES
%----------------------------------------- -----------------------------------------------
\newcommand{\mystrut}{\rule[-.3\baselineskip]{0pt}{\baselineskip}}

%----------------------------------------------------------------------------------------
% CUSTOM LOREM IPSUM
%----------------------------------------------------------------------------------------
\newcommand{\lorem}
{Lorem ipsum dolor sit amet, consectetur adipiscing elit. Donec a diam lectus.}



%============================================================================%
%
%
%
%    DOCUMENT CONTENT
%
%
%
%============================================================================%
\begin{document}

%use our custom fancy header definitions
\pagestyle{fancy}

%----------------------------------------------------------------------------------------
%    HEADER IMAGE
%----------------------------------------------------------------------------------------

%\begin{figure}[H]
%\begin{flushright}
%    \includegraphics[trim= 320 130 460 210,clip,width=0.2\linewidth]{myfoto.jpg}    %trimming relative to image size!
%\end{flushright}
%\end{figure}

%---------------------------------------------------------------------------------------
%    TITLE HEADLINE
%----------------------------------------------------------------------------------------
\vspace{-8pt}
\begin{center}
    \HUGE \textsc{Jason Fevang} \textcolor{sectcol}{\rule[-1mm]{1mm}{0.9cm}} \textsc{Resume}\\[2pt]
    \small Software Engineer
\end{center}

\vspace{6pt}

%---------------------------------------------------------------------------------------
%    META SECTION
%----------------------------------------------------------------------------------------
\metasection{Vancouver, Canada}{\textbf{Software:} C++11/20, Python, FreeRTOS}
\metasection{jasonfevang@gmail.com}{\textbf{Hardware:} ESP32/8266, STM8, PWM}
\metasection{github.com/jasonfevang}{\textbf{Protocols:} UART, SPI, I2C, UDP/HTTP (lwIP Sockets)}
\metasection{604-401-4643}{\textbf{Tools:} Oscilloscopes, Logic Analyzers, Agentic AI, Unit Testing (Pytest/GTest)}
\vspace{-2pt}
\textcolor{softcol}{\hrule}
\vspace{6pt}

\normalsize

%---------------------------------------------------------------------------------------
%    SUMMARY (optional)
%----------------------------------------------------------------------------------------
\vspace{-6pt}
\cvsection{Summary}
I'm a computer engineering graduate with seven years of work experience in embedded and enterprise development with C++ and Python. I've developed firmware for professional LED lighting products on FreeRTOS+ESP32 and baremetal STM8. I'm currently implementing data formats focused on memory safety, security, and performance at Safe Software.

%============================================================================%
%
%    CV SECTIONS AND EVENTS (MAIN CONTENT)
%
%============================================================================%

%---------------------------------------------------------------------------------------
%    EXPERIENCE
%----------------------------------------------------------------------------------------
\cvsection{Work Experience}

\cvevent{2022 - Current (3.5 years)}{C++ Software Developer}{Safe Software}{
    {Developed high-performance features for a cross-platform enterprise application using C++20, prioritizing memory safety, system reliability and protection against security vulnerabilities.},
    {Managed the full software development lifecycle, including code reviews and CI/CD pipelines, while maintaining detailed documentation for all tasks to ensure technical clarity and traceability.},
    {Implemented an AI-augmented development workflow using agentic tools (Claude-Code, Copilot) and authored a suite of shell scripts to automate repetitive development tasks, increasing team productivity.}
}

\cvevent{2019 - 2022 (3 years)}{Embedded Software Developer}{Better Way Lighting}{
    {Performed board bring-up for new hardware, validating and troubleshooting with oscilloscopes and logic analyzers.},
    {Integrated hardware components including displays(SSD1306, ILI9341), sensors, and communication modules, utilizing UART, SPI, and I2C protocols.},
    {Collaborated with electrical engineers and designers to develop firmware with real-time constraints for lighting products used in major film productions.},
    {Oversaw work terms for four software co-op students, performing interviews, performance reviews, mentorship and project management.}
}

%---------------------------------------------------------------------------------------
%    PROJECTS SECTION
%--------------------------------------------------------------------------------------
\vspace{-10pt}
\cvsection{Selected Projects}

% \cvevent{2025}{Python PDF Writer}{Safe Software}{
%     {Designed, implemented and \href{https://docs.safe.com/fme/html/FME-Form-Documentation/FME-ReadersWriters/pkg-documentpdf/Feature_Representation.htm}{documented} an output format extending FPDF2 for creating text-based PDF reports that has been downloaded 1300 times in its first nine months.},
%     {Wrote 215 unit tests covering a broad feature-set including multi-column layouts, custom fonts, images, bullet points, watermarks and much more.}
% }

\cvevent{2025}{Optimize GIS PointCloud Network Access}{Safe Software}{
    {Built a high-performance C++ implementation of Cloud Optimized Point Cloud (\href{https://copc.io}{copc.io}), using an octree to spatially index datasets upwards of 1B points.},
    {Optimized remote datasets with HTTP/FTP range requests and over-fetching, avoiding full-file downloads.}
}

\cvevent{2022}{LimeLite LED Panel Firmware - \href{https://betterwaylighting.com/limelite}{link}}{Better Way Lighting}{
    {Architected ESP32 firmware on FreeRTOS to handle real-time lighting data over UDP (\href{https://tsp.esta.org/tsp/documents/docs/E1-31-2016.pdf}{sACN}/\href{https://art-net.org.uk/downloads/art-net.pdf}{Art-Net}), physical DMX512, and an external networking device via SPI (\href{https://docs.lumenrad.io/timo/spi-interface/}{CRMX}), writing all libraries.},
    {Implemented closed-loop thermal control with thermistors and PWM-driven fans, including ramp-up logic.},
    % {Developed a responsive user interface, managed via FreeRTOS tasks to ensure high-performance UI interaction alongside background networking.},
    {Engineered a high-fidelity LED driving system using variable PWM with gamma correction and look-up tables for color mixing and panel-specific calibration to achieve \textgreater 95 CRI white light across the spectrum.}
}

\cvevent{2022}{STM8 Decoder Firmware and Bootloader - \href{https://betterwaylighting.com/obd-decoder}{link}}{Better Way Lighting}{
    {Developed custom firmware in C for STM8 microcontrollers, implementing DMX512 receive via UART and driving character displays and LEDs through low-level GPIO and TIM timer peripherals.},
    % {Reverse-engineered and wrote custom firmware in C for six varieties of DMX512 decoders, with over 1000 devices programmed.},
    {Wrote a custom bootloader and programmer (ESP32) supporting custom commands for field updates.}
    % {Automated all configuration, build, and flashing tooling using makefiles and python scripts for development with the COSMIC STM8 Toolchain.}
}

% \cvevent{2021}{C++ GUI Framework for Monochrome Displays}{Better Way Lighting}{
%     {Designed and built CardsGUI, used by the development team to build GUI's now running on multiple product lines.},
%     {Implemented in object-oriented C++ on top of U8g2, using FreeRTOS primitives on the ESP32.}
% }

% \cvevent{2021}{Handheld Dmx Multitool}{Better Way Lighting}{
%     {Developed a full-color UI using LVGL, abstracted such that UI development occurred on desktop, then was cross-compiled to the ESP32.},
%     {Led an agile-based scrum team of two members to coordinate development efforts.},
%     {Architected application software in C++, supporting hardware features such as SD card use, DMX512, Wi-Fi, battery management and more.}
% }

%---------------------------------------------------------------------------------------
%    EDUCATION SECTION
%--------------------------------------------------------------------------------------
\cvsection{Education}

\cvevent{2015 - 2019}{BASc Computer Engineering}{Simon Fraser University}{
    {Bachelor of Applied Science, Computer Engineering (With Distinction), 3.92 GPA}
}

%--------------------------------------------------------------------------------------------------
%    ARTIFICIAL FOOTER (fancy footer cannot exceed linewidth) 
%--------------------------------------------------------------------------------------------------

\null
\vspace*{\fill}
\hspace{-0.25\linewidth}\colorbox{white}{\makebox[1.5\linewidth][c]{\mystrut  \textnormal{\textcolor{sectcol}{jasonfevang@gmail.com} $\cdot$ \textcolor{sectcol}{github.com/jasonfevang}}}}

%============================================================================%
%
%
%
%    DOCUMENT END
%
%
%
%============================================================================%
\end{document}
